% 08_conclusion.tex — 结论与演进路径

\section{结论与演进路径}

\subsection{主要结论}

\subsubsection{1. LEO 并非太空算力的最优轨道}

直觉上 LEO 应该最便宜（最低发射成本、最低延迟），但综合评估表明：
\begin{itemize}
    \item GEO 综合成本指数 0.83，比 LEO 优 17\%——15 年已验证寿命 + 在轨服务 + 小光伏面积 + 简单热设计，在"可维护性"和"全生命周期"维度上对 LEO 有根本性优势
    \item LEO 的 3--5 年硬件寿命导致的替换成本可能超过 GEO 的高发射成本
    \item 地月 L4/L5（0.94）和地日 L1/L2（0.99）同样在综合指数上不劣于 LEO
\end{itemize}

\subsubsection{2. TID 悖论：深空辐射挑战的本质是 SEE}

深空（地日 L1/L2）在适度屏蔽后的年 TID 仅为 2--5~\kradyr——\textit{低于} 穿越 SAA 的 LEO 轨道（10--30~\kradyr）。但这是误导性的"优势"：SEU 率比 LEO 高 100--1000 倍，且 SEE 防护（TMR + EDAC + scrubbing）是持续性的面积/功耗惩罚，无法通过一次性工艺选择消除。"猝死"（SEE）比"慢死"（TID）更难应对。

\subsubsection{3. MEO 和 HEO 不推荐部署太空算力}

MEO（SEU 率 $\sim 10^{5}$x LEO）和 HEO（TID 50--200~\kradyr + SEU $\sim 10^{5}$--$10^{6}$x LEO）的辐射环境使防护开销可能"吃掉"大部分有效算力。两者均缺乏独特的应用价值主张。

\subsubsection{4. 通信延迟是不可消除的硬约束}

地日 L1/L2 的 10 秒 RTT（1000--2000x LEO）排除交互式推理，仅适用于"批处理式"计算。地月 L4/L5 的 2.6 秒 RTT 是折中方案。GEO 的 480--600 ms RTT 可通过本地处理规避。

\subsection{非 LEO 太空算力优先级排序}

\newcommand{\stars}[1]{\textbf{\textcolor{accent}{#1}}}

\begin{table}[H]
\centering
\caption{非 LEO 太空算力部署优先级排序}
\label{tab:priority}
\small
\begin{tabular}{@{}cllp{5.5cm}@{}}
\toprule
\textbf{梯队} & \textbf{轨道} & \textbf{时间窗口} & \textbf{最可能场景} \\
\midrule
Tier 1 & GEO        & 2028--2035   & 通信卫星"搭载"AI推理载荷，区域实时推理广播 \\
Tier 2 & 地月 L4/L5 & 2032--2040+  & 深空通信/导航骨干计算节点（LunaNet集成）\\
Tier 3 & 地日 L1/L2 & 2035--2040+  & L2量子-经典协同计算（唯一独特物理优势）\\
Tier 4 & MEO        & 2035+ (特定) & GNSS星座原位AI增强（仅边缘应用）\\
Tier 5 & HEO        & 不推荐       & 无可行场景 \\
\bottomrule
\end{tabular}
\end{table}

\subsection{演进时间线}

\begin{description}
    \item[2026--2028\quad 原型验证期] LEO 是唯一有硬件验证计划的轨道（AI1/Exlumina 原型卫星）。DARPA RSGS 首次 GEO 在轨服务演示——这是 GEO 算力的关键使能事件。非 LEO 算力为零实际部署。
    \item[2028--2032\quad GEO 先行部署期] 首批通信卫星"算力载荷"入轨（10--50 kW 级别 AI 推理）。军事/情报机构最可能为首批客户。在轨服务（RSGS/MRV）开始商业运营。L4/L5 方向：Gateway NRHO 就位，LunaNet 规划中。
    \item[2032--2037\quad 多轨道扩展期] GEO 算力载荷规模扩大至 100 kW+ 级别。若 Artemis 月球基地常态化 $\rightarrow$ L4/L5 深空中继计算节点部署。L2 量子协同计算概念验证（若超导量子-空间集成技术突破）。
    \item[2037--2040\quad 多轨道生态形成期] GEO 成熟商业市场（通信卫星运营商标配 AI 载荷）。L4/L5 月球-深空互联网骨干计算节点。L1/L2 可能首次专用算力任务。
\end{description}

\subsection{关键先行事件}

\begin{table}[H]
\centering
\caption{关键先行事件及其重要性}
\label{tab:key-events}
\small
\begin{tabular}{@{}clcl@{}}
\toprule
\textbf{\#} & \textbf{事件} & \textbf{预计时间} & \textbf{影响轨道} \\
\midrule
1 & DARPA RSGS 首次 GEO 在轨服务成功 & 2027--2028 & GEO \hspace{4em} \\
2 & AI1 卫星在轨散热/可靠性数据 & 2027--2028 & LEO $\rightarrow$ 所有 \\
3 & Starship 在轨加注常态化 & 2028--2029 & 所有高轨道 \\
4 & SpaceX Terafab D3 轨道专用芯片量产 & 2028--2030 & 所有 \\
5 & Artemis 月球基地常态化运营 & 2030--2032 & L4/L5 \\
6 & 超导量子计算空间环境集成验证 & 2032--2035 & L2 \\
7 & 太空光伏 \$/W 降至 $<$\$100 & 2030--2035 & 所有 \\
\bottomrule
\end{tabular}
\end{table}

\subsection{最终判断}

非 LEO 太空算力最可能演进路径为：

\begin{center}
\Large
\textbf{GEO 先行（2028--2035）$\rightarrow$ L4/L5 跟随（2032--2040）$\rightarrow$ L2 利基应用（2035+）}
\end{center}

\bigskip

GEO 是最可能率先实现经济可行的非 LEO 轨道——不是因为发射成本低，而是因为 15 年寿命验证 + 在轨服务降低了全生命周期成本。地月 L4/L5 是最有战略价值的远期资产，但需要 Artemis/LunaNet 生态成熟。L2 量子协同是唯一值得 L2 成本溢价的利基场景。MEO 和 HEO 不推荐任何太空算力部署。

\textbf{降低硬件抗辐射溢价——而非降低发射成本——才是使任何轨道太空算力经济可行的最关键杠杆。}
